\section{Integrated stress}\label{sec:integrated-stress}

% Mathematical job:
%   - Present the six stress families as anti-overclaim guardrails, not as
%     proofs of the main theorem:
%       1. dependency stress (R010, lower bounds, upper bound,
%          decomposition present in order)
%       2. circularity stress (no circular domain, no by-definition
%          no-seventh-role, no decomposition assuming exact-six)
%       3. channel-vs-activation stress (inactive != active)
%       4. role-alignment stress (no generic transition->P3 etc.)
%       5. upper-bound residual stress (probability/causality/agency do
%          not reopen; no unclassified residual; no genuine seventh)
%       6. overclaim stress (no full algebra, no global independence,
%          no unique decomposition, no empirical realization, no
%          unrestricted all-emergence)
%   - Summarize with tables; state promotion-readiness.
%
% Governing sources:
%   workspace/meta/lens_lab/campaigns/paper3_exact_six_exhaustiveness/integrated_stress/exact_six_integrated_stress_suite.md
%   workspace/meta/lens_lab/aristotle/results/paper3/paper3_exact_six_integrated_stress_H/paper3_exact_six_integrated_stress_H_result_shape.md
%   workspace/meta/lens_lab/papers/paper3/drafting_support/figure_table_plan.md
%     (table_5_integrated_stress)
%
% Overclaim risks:
%   - Presenting stress cases as if they were the proof of the theorem
%     rather than stability checks on the theorem stack.
%   - Converting stress pass into universal validity.


With the theorem chain of Sections~\ref{sec:lower-bounds},
\ref{sec:upper-bound-classification}, and
\ref{sec:decomposition-exhaustion} in place, we record six stress
families that act as consistency checks against overclaims. These are
not proofs of the scoped exact-six target; they verify that the
theorem chain does not silently admit overclaims. Their joint
satisfaction is the last requirement before the terminal statement in
Section~\ref{sec:scoped-exact-six}.

\subsection{Dependency stress}\label{subsec:dependency-stress}

\begin{proposition}[Dependency stress]\label{prop:dependency-stress}
The scoped exact-six target is dependency-coherent only when the theorem
chain is assembled in the order
\[
\text{\FATCD{} domain} \;\rightarrow\; \text{lower bounds} \;\rightarrow\;
\text{upper-bound classification} \;\rightarrow\;
\text{decomposition/exhaustion}.
\]
Accordingly, the following are inadmissible:
\begin{itemize}[topsep=2pt,itemsep=1pt]
  \item asserting the target without the \FATCD{} domain of
    Section~\ref{sec:fatcd-domain};
  \item asserting the target without the lower-bound theorem of
    Section~\ref{sec:lower-bounds};
  \item asserting the target without the upper-bound classification
    theorem of Section~\ref{sec:upper-bound-classification};
  \item asserting the target without the decomposition and exhaustion
    theorems of Section~\ref{sec:decomposition-exhaustion};
  \item promoting the assembled target before the integrated stress
    check of Theorem~\ref{thm:integrated-stress}.
\end{itemize}
\end{proposition}

\subsection{Circularity stress}\label{subsec:circularity-stress}

\begin{proposition}[Circularity stress]\label{prop:circularity-stress}
Three circularity patterns are inadmissible under the theorem chain:
\begin{itemize}[topsep=2pt,itemsep=1pt]
  \item redefining \FATCD{} so that it already contains exactly six role
    slots by construction;
  \item declaring the absence of a seventh role by definition rather
    than by classification under Definition~\ref{def:residual-scheme};
  \item justifying decomposition existence by assuming scoped
    exact-six.
\end{itemize}
The non-circularity schema of Proposition~\ref{prop:non-circularity} and
the constructive proof of Theorem~\ref{thm:decomposition-existence}
block each pattern.
\end{proposition}

\subsection{Channel-versus-activation stress}\label{subsec:channel-activation-stress}

\begin{proposition}[Channel-versus-activation stress]\label{prop:channel-activation-stress}
The following conflations are inadmissible:
\begin{itemize}[topsep=2pt,itemsep=1pt]
  \item requiring that every \(D \in \FATCD{}\) activates all six role
    channels;
  \item treating an inactive channel of status \(\belowthr\) as an
    active witness for its role;
  \item treating an inactive channel of status \(\collapsedbc\) as an
    active witness;
  \item treating an inactive channel of status \(\absentrec\) as an
    active witness.
\end{itemize}
Proposition~\ref{prop:channel-vs-activation} enforces each distinction.
The same inactive-status guardrail also applies to
\(\outsidescope\).
\end{proposition}

\subsection{Role-alignment stress}\label{subsec:role-alignment-stress}

\begin{proposition}[Role-alignment stress]\label{prop:role-alignment-stress}
The three forbidden assignments of
Remark~\ref{rem:three-repaired-alignments} are inadmissible: generic
transition data does not project to \Pthree{}, generic path or
derivation data does not project to \Pfive{}, and generic semantic or
interpretation data does not project to \Pone{}. The alignment rules
of Proposition~\ref{prop:alignment-rules} enforce each prohibition.
\end{proposition}

\subsection{Upper-bound residual stress}\label{subsec:upper-bound-residual-stress}

\begin{proposition}[Upper-bound residual stress]\label{prop:residual-stress}
Under Theorem~\ref{thm:upper-bound-classification} and within the covered
scope, the following regressions are inadmissible:
\begin{itemize}[topsep=2pt,itemsep=1pt]
  \item probability or uncertainty features reopening as unresolved
    threats;
  \item residual-causality features reopening as unresolved threats;
  \item residual-agency features reopening as unresolved threats;
  \item any covered residual feature remaining unclassified under
    Definition~\ref{def:residual-scheme};
  \item any covered residual feature requiring a genuine seventh
    irreducible typed closure role.
\end{itemize}
Propositions~\ref{prop:probability-closure}--\ref{prop:agency-closure}
settle the three threat closures, and
Theorem~\ref{thm:upper-bound-classification} delivers the emptiness of
categories (x) and (xi).
\end{proposition}

\subsection{Overclaim stress}\label{subsec:overclaim-stress}

\begin{proposition}[Overclaim stress]\label{prop:overclaim-stress}
The following restatements of the scoped exact-six target are
inadmissible:
\begin{itemize}[topsep=2pt,itemsep=1pt]
  \item restatement as a full six-primitives algebra;
  \item restatement as a global primitive-independence theorem;
  \item restatement as unique or canonical decomposition;
  \item restatement as empirical realization;
  \item restatement as a claim about all emergence phenomena, or any
    domain broader than \FATCD{}.
\end{itemize}
Remark~\ref{rem:atlas-global-nonclaims},
Proposition~\ref{prop:nonclaim-closure}, and
Proposition~\ref{prop:non-uniqueness} carry the corresponding nonclaims
forward through the theorem chain. In particular, nonclaim closure
survives the exact-six assembly.
\end{proposition}

\subsection{The integrated stress theorem}\label{subsec:integrated-stress-theorem}

\begin{theorem}[Integrated stress]\label{thm:integrated-stress}
The theorem chain of Sections~\ref{sec:fatcd-domain},
\ref{sec:lower-bounds}, \ref{sec:upper-bound-classification}, and
\ref{sec:decomposition-exhaustion} satisfies each of the six stress
families stated in
Propositions~\ref{prop:dependency-stress}--\ref{prop:overclaim-stress}.
Dependency, circularity, channel-versus-activation, role-alignment,
upper-bound residual, and overclaim regressions are all blocked, and
the assembled target is ready for the terminal statement of
Section~\ref{sec:scoped-exact-six}.
\end{theorem}

\begin{proof}
Each family is discharged by its corresponding proposition.
Dependency stress follows from the ordered presence of the \FATCD{}
domain (Section~\ref{sec:fatcd-domain}), lower bounds
(Section~\ref{sec:lower-bounds}), upper-bound classification
(Section~\ref{sec:upper-bound-classification}), and
decomposition/exhaustion (Section~\ref{sec:decomposition-exhaustion}).
Circularity stress follows from
Proposition~\ref{prop:non-circularity} and the constructive proof of
Theorem~\ref{thm:decomposition-existence}. Channel-versus-activation
stress follows from Proposition~\ref{prop:channel-vs-activation} and
the five-status taxonomy of Definition~\ref{def:channel-status}.
Role-alignment stress follows from
Proposition~\ref{prop:alignment-rules} and
Remark~\ref{rem:three-repaired-alignments}. Upper-bound residual
stress follows from
Propositions~\ref{prop:probability-closure}--\ref{prop:agency-closure}
and Theorem~\ref{thm:upper-bound-classification}. Overclaim stress
follows from Remark~\ref{rem:atlas-global-nonclaims},
Proposition~\ref{prop:nonclaim-closure}, and
Proposition~\ref{prop:non-uniqueness}. Together: dependencies cohere,
the channel/activation distinction holds, nonclaim closure holds, and
the target is eligible for the terminal statement.
\end{proof}
